% ==============================================================================
% CAPITOLO 12 - FONDAMENTI DI TERMODINAMICA, TEMPERATURA E CALORE
% ==============================================================================
\chapter{Fondamenti di Termodinamica, Temperatura e Calore}
\label{chap:termodinamica_fondamenti}

\begin{tcolorbox}[enhanced, breakable, colback=navyblue!4!white, colframe=navyblue!80!black,
    coltitle=white, fonttitle=\bfseries\sffamily, title={Quadro Sinottico del Capitolo 12 (Corrispondenza: Lezioni 32, 33 e 34)},
    sharp corners, rounded corners=south, boxrule=1pt]
\textbf{Collocazione Modulare:} Parte IV -- Termodinamica\\
\textbf{Trascrizioni di Riferimento:} \texttt{trascrizioni/Lezione\_32\_...md}, \texttt{Lezione\_33\_...md}, \texttt{Lezione\_34\_...md}\\
\textbf{Manuale di Riferimento:} Focardi, Massa, Ugolini -- \emph{Fisica Generale: Meccanica e Termodinamica} (Capitoli 13 e 14)\\
\textbf{Obiettivi Didattici e Competenze:}\par
\begin{itemize}
    \item Definire i sistemi termodinamici (aperti, chiusi, isolati), pareti (adiabatiche, diatermiche), variabili intensive ed estensive.
    \item Formulare il Principio Zero della Termodinamica e definire rigorosamente la temperatura empirica e la scala termodinamica assoluta Kelvin.
    \item Analizzare la dilatazione termica e il calore $Q$ come energia microscopica disordinata in transito (equivalente meccanico di Joule).
    \item Calcolare capacità termiche, calori specifici e transizioni di fase con calore latente in problemi di calorimetria.
    \item Modellizzare i tre meccanismi di trasmissione del calore: Conduzione (Legge di Fourier e resistenze termiche), Convezione (Newton) e Irraggiamento (Legge di Stefan-Boltzmann).
\end{itemize}
\end{tcolorbox}

\vspace{0.5em}

% ------------------------------------------------------------------------------
\section{Sistemi Termodinamici, Variabili di Stato ed Equilibrio}
\label{sec:sistemi_termodinamici}
% ------------------------------------------------------------------------------

La \textbf{termodinamica} è la teoria macroscopica e fenomenologica che studia le trasformazioni dell'energia sotto forma di lavoro e calore e le proprietà della materia all'equilibrio.

\begin{definizione}{Classificazione dei Sistemi Termodinamici}{def:classificazione_sistemi}
Un \textbf{sistema termodinamico} è una porzione finita di materia separata dal resto dell'universo (ambiente circostante) da un contorno reale o ideale:
\begin{itemize}
    \item \textbf{Aperto}: scambia sia materia che energia (sotto forma di calore e lavoro) con l'ambiente.
    \item \textbf{Chiuso}: scambia energia ma non materia con l'ambiente.
    \item \textbf{Isolato}: non scambia né materia né energia di alcun tipo con l'ambiente ($\Delta M = 0, Q = 0, W = 0$).
\end{itemize}
\end{definizione}

\begin{definizione}{Variabili Intensive ed Estensive}{def:variabili_stato}
\begin{itemize}
    \item \textbf{Intensive}: indipendenti dalle dimensioni e dall'estensione spaziale del sistema (es. pressione $p$, temperatura $T$, densità $\rho$, potenziale chimico $\mu$).
    \item \textbf{Estensive}: proporzionali alla massa o al volume totale del sistema (es. volume $V$, massa $M$, energia interna $U$, entropia $S$, entalpia $H$).
\end{itemize}
\end{definizione}

% ------------------------------------------------------------------------------
\section{Il Principio Zero della Termodinamica e la Temperatura}
\label{sec:principio_zero}
% ------------------------------------------------------------------------------

\begin{legge}{Principio Zero della Termodinamica}{legge:principio_zero}
Se due corpi $A$ e $B$ sono ciascuno in equilibrio termico con un terzo corpo $C$ (termometro), allora $A$ e $B$ sono in reciproco equilibrio termico tra loro:
\begin{equation}
    (A \sim C) \land (B \sim C) \implies A \sim B
\end{equation}
\end{legge}

Il Principio Zero garantisce l'esistenza di una funzione di stato scalare universale, la \textbf{temperatura} $T$: due sistemi in contatto termico sono in equilibrio se e solo se possiedono la medesima temperatura ($T_A = T_B$).

\subsection{Scale Termometriche e Temperatura Assoluta}
Nel Sistema Internazionale l'unità di misura della temperatura assoluta è il \textbf{Kelvin} ($\si{\kelvin}$), fissato stabilendo lo zero assoluto a $\SI{0}{\kelvin}$ (temperatura in cui l'energia cinetica molecolare classica si annulla):
\begin{equation}
    T(\si{\kelvin}) = T(^\circ\text{C}) + 273{,}15
\end{equation}

\subsection{Dilatazione Termica dei Solidi e dei Liquidi}
Al variare della temperatura $\Delta T = T - T_0$:
\begin{itemize}
    \item \textbf{Dilatazione Lineare}: $\Delta L = \alpha L_0 \Delta T \implies L(T) = L_0(1 + \alpha\Delta T)$ (con $\alpha$ coeff. di dilatazione lineare, $\si{\per\kelvin}$).
    \item \textbf{Dilatazione Volumica}: $\Delta V = \beta V_0 \Delta T \approx 3\alpha V_0 \Delta T$ (per solidi isotropi $\beta = 3\alpha$).
\end{itemize}

% ------------------------------------------------------------------------------
\section{Calore, Capacità Termica e Calorimetria}
\label{sec:calore_calorimetria}
% ------------------------------------------------------------------------------

\begin{definizione}{Calore}{def:calore}
Il \textbf{calore} $Q$ è l'energia termica microscopica trasferita spontaneamente attraverso il contorno di un sistema per effetto esclusivo di una differenza di temperatura $\Delta T$ con l'ambiente ($T_{\text{corpo}} \neq T_{\text{amb}}$).
Nel SI si misura in **Joule** ($\si{\joule}$). L'\textbf{equivalente meccanico del calore} di Joule stabilisce:
\begin{equation}
    \SI{1}{\calorie} = \SI{4.1868}{\joule}
\end{equation}
\end{definizione}

\subsection{Capacità Termica e Calore Specifico}
\begin{itemize}
    \item \textbf{Capacità Termica} $C$: $C \coloneqq \dv{Q}{T} \quad (\si{\joule\per\kelvin})$.
    \item \textbf{Calore Specifico Massico} $c$: quantità di calore necessaria per innalzare di $\SI{1}{\kelvin}$ la temperatura di $\SI{1}{\kilogram}$ di sostanza:
    \begin{equation}
        c \coloneqq \frac{C}{m} = \frac{1}{m}\dv{Q}{T} \implies \mathbf{Q = m c \Delta T = m c (T_f - T_i)}
        \label{eq:calore_sensibile}
    \end{equation}
    Per l'acqua liquida: $c_{\text{acqua}} \approx \SI{4186}{\joule\per\kilogram\per\kelvin} = \SI{1}{\calorie\per\gram\per\celsius}$.
\end{itemize}

\subsection{Cambiamenti di Fase e Calore Latente}
Durante una transizione di fase isobara e isoterma (es. fusione o ebollizione), la temperatura rimane rigorosamente costante. Il calore scambiato vale:
\begin{equation}
    \mathbf{Q = \pm m L}
    \label{eq:calore_latente}
\end{equation}
dove $L$ è il \textbf{calore latente} della sostanza ($\si{\joule\per\kilogram}$, con segno positivo per fusione/evaporazione, negativo per solidificazione/condensazione).

\subsection{Equilibrio Calorimetrico (Principio di Conservazione dell'Energia)}
In un calorimetro adiabatico ideale ($Q_{\text{scambiato con esterno}} = 0$):
\begin{equation}
    \sum_{i=1}^N Q_i = 0 \implies \sum_{i=1}^N m_i c_i (T_{\text{eq}} - T_i) = 0 \implies \mathbf{T_{\text{eq}} = \frac{\sum m_i c_i T_i}{\sum m_i c_i}}
    \label{eq:temperatura_equilibrio}
\end{equation}

% ------------------------------------------------------------------------------
\section{I Tre Meccanismi di Trasmissione del Calore}
\label{sec:trasmissione_calore}
% ------------------------------------------------------------------------------

\subsection{Conduzione Termica (Legge di Fourier)}
Trasferimento di calore a livello molecolare/reticolare nei solidi senza moto macroscopico di materia:
\begin{equation}
    \mathbf{\dot{Q} \coloneqq \dv{Q}{t} = -k A \dv{T}{x} \xrightarrow{\text{parete piana}} \dot{Q} = k A \frac{T_1 - T_2}{L} = \frac{T_1 - T_2}{R_{\text{th}}}}
    \label{eq:fourier}
\end{equation}
dove $k$ è la \textbf{conducibilità termica} ($\si{\watt\per\metre\per\kelvin}$), $A$ è l'area, $L$ lo spessore e $R_{\text{th}} \coloneqq \frac{L}{k A}$ è la **resistenza termica conduttiva** ($\si{\kelvin\per\watt}$).

\subsection{Convezione Termica (Legge di Newton del Raffreddamento)}
Scambio termico tra una superficie solida a temperatura $T_s$ e un fluido in movimento a temperatura asintotica $T_\infty$:
\begin{equation}
    \mathbf{\dot{Q} = h A (T_s - T_\infty)}
    \label{eq:convezione}
\end{equation}
dove $h$ è il coefficiente di scambio termico convettivo ($\si{\watt\per\metre\squared\per\kelvin}$).

\subsection{Irraggiamento Termico (Legge di Stefan-Boltzmann)}
Emissione di radiazione elettromagnetica da parte di qualsiasi corpo a temperatura $T > \SI{0}{\kelvin}$:
\begin{equation}
    \mathbf{P_{\text{em}} = \varepsilon \sigma A T^4}
    \label{eq:stefan_boltzmann}
\end{equation}
dove $\sigma = \SI{5.670374e-8}{\watt\per\metre\squared\per\kelvin\tothe{4}}$ è la costante di Stefan-Boltzmann ed $\varepsilon \in [0, 1]$ è l'\textbf{emissività} della superficie ($\varepsilon=1$ per il corpo nero ideale). Per un corpo immerso in un ambiente a temperatura $T_{\text{amb}}$:
\begin{equation}
    P_{\text{net}} = \varepsilon \sigma A (T^4 - T_{\text{amb}}^4)
\end{equation}
